\section{Limitations}

The all-size theorems concern two exact algebraic grammars and frozen structural objectives.  They do not establish physical runtime, fault-tolerant resource savings, hardware advantage or production-scale compiler behavior.  The six-term result is likewise restricted to its admitted batch family and cost decomposition.

The state-preparation conclusions are finite.  Complete state graphs are used only through $n=3$, while transfer is tested on one frozen 120-state $n=4$ panel.  The 127-feature representation is sign aware and permutation covariant, but it is near-injective in domain and has negligible parity-split coverage on the panel.  Other representations, gate sets, costs or measurement-assisted preparation may behave differently.

The objective inequalities describe a sufficient region for one support-one proof.  They are not proved necessary.  A corrected frozen-domain search found no strict outside witness, but a witness in a broader domain remains possible.

The rank-minus-support observation uses only two measured models and changes under a different rewrite alignment.  It should be treated as a diagnostic coincidence until a third model discriminates between competing explanations.

Independent scripts and verifier implementations remain under one project and one custody chain.  They check consistency and structural separation, but they are not independent external replication or external proof review.  Public archival identity and third-party verification remain separate publication tasks.
